\section{Grobe Flow-Analyse}
\label{sec:flow}

Dieser Teil rechnet. Er beantwortet zwei Fragen, und sie schlie{\ss}en
einander aus: ob das ver\"offentlichte Wachstumsprogramm aus eigener Kraft
finanzierbar ist (Abschnitt~\ref{sec:kapazitaet}), und was ein Anleger
verdient, der heute einsteigt (Abschnitt~\ref{sec:anleger}). Dieselben Mittel
k\"onnen nur einmal verwendet werden. Die erste Rechnung leitet jeden
\"Uberschuss in die Investition, die zweite an die Eigent\"umer; sie
nebeneinander zu lesen und beide Ertr\"age zu addieren hie{\ss}e, dasselbe
Geld zweimal auszugeben.

Eine Entschuldungsrechnung, wie sie bei einem verschuldeten Unternehmen an
erster Stelle st\"unde, entf\"allt: Alamos h\"alt nach
Abschnitt~\ref{sec:bilanz} mehr Kasse als Finanzschulden.

\subsection{Annahmen}
\label{sec:annahmen}

Die Betr\"age dieser Rechnung sind belegt, ihre Mechanik ist es nicht. Jede
nicht aus den Quellen ableitbare Festlegung steht deshalb hier, vor der
Rechnung, die sie verbraucht.

\annahme{\textbf{Preisszenarien.} Gerechnet werden drei Goldpreise, und keiner
davon ist gesch\"atzt: Es sind die drei vom Unternehmen selbst berichteten
\emph{durchschnittlich erzielten} Preise -- \JeUnze{\GoldpreisVJ} f\"ur 2024,
\JeUnze{\Goldpreis} f\"ur 2025 und \JeUnze{\HJGoldpreis} f\"ur das erste
Halbjahr 2026 -- jeweils konstant fortgeschrieben. Damit muss keine
Preisannahme erfunden werden. Dass der Preis und nicht das Mengenwachstum die
tragende Gr\"o{\ss}e ist, folgt aus Abschnitt~\ref{sec:margejeunze}.}

\annahme{\textbf{Menge und St\"uckkosten.} Sie folgen dem vom Unternehmen am
\Datum{4.\,Februar 2026} ver\"offentlichten Ausblick: bis 2027 die Mitte der
Prognose 2026 mit \TsdUnzen{\FlowMengeSechs} bei AISC von
\JeUnze{\FlowAiscSechs}, ab 2028 die Mitte des Dreijahresausblicks mit
\TsdUnzen{\FlowMengeAcht} bei AISC von \JeUnze{\FlowAiscAcht}, das sind
\Pct{\ProgAchtAISCRueckgang} unter dem Stand von 2025. Verwendet wird die im
Zwischenbericht zum \Datum{30.\,Juni 2026} \emph{angehobene} Kostenprognose;
die urspr\"ungliche w\"are die g\"unstigere und damit die unehrlichere Wahl.
F\"ur 2027 ver\"offentlicht das Unternehmen keine eigene Zahl; das Jahr wird
auf dem Stand von 2026 gehalten und nicht zwischen 2026 und 2028
interpoliert, weil eine Interpolation eine erfundene Zwischenstufe w\"are.
Die Zur\"uckhaltung f\"allt zulasten der Rechnung.}

\annahme{\textbf{Abgrenzung des Cash-Flows.} Die Rechnung setzt auf der
Kostenkennzahl auf, die das Unternehmen selbst ver\"offentlicht. Die All-in
Sustaining Costs enthalten das Erhaltungskapital bereits, nicht aber das
Wachstumskapital und nicht die aktivierte Exploration. Der verteilbare
Cash-Flow lautet deshalb
\[
  \text{Menge} \times (\text{Goldpreis} - \text{AISC})
  - \text{Wachstumskapital} - \text{Exploration} - \text{Steuer}
\]
und nicht \enquote{operativer Cash-Flow abz\"uglich Investitionen}: Die zweite
Form lie{\ss}e sich nicht nach dem Goldpreis aufl\"osen, und der Goldpreis ist
die Gr\"o{\ss}e, die dieses Gesch\"aft bewegt. Abschnitt~\ref{sec:eichung}
h\"alt das Modell gegen das Ist-Jahr 2025.}

\annahme{\textbf{Dauer des Bauprogramms.} Das Wachstumskapital wird mit
\MioUSDexakt{\FlowWachstumskapital} je Jahr angesetzt, der Mitte der Prognose
2026, und endet mit dem Jahr \num{\FlowBauende} -- dem Jahr, f\"ur das das
Unternehmen die erste Produktion von Lynn Lake angek\"undigt hat. Diese
Annahme entscheidet das Ergebnis mit und wird in
Tabelle~\ref{tab:flowsensitiv} in beide Richtungen durchgerechnet. Sie mit
dem Spitzenbetrag dauerhaft fortzuschreiben und zugleich die Produktion auf
dem Stand von 2026 zu halten w\"are kein vorsichtiger Ansatz, sondern ein
falscher: Die Rechnung bezahlte dann f\"unfzehn Jahre lang eine Erweiterung
und erhielte sie nie.}

\annahme{\textbf{Steuern.} Angesetzt wird der im Abschluss genannte
gesetzliche Satz von \Pct{\Steuersatz}\quelleFS{2025}{34}\quelleMerken{s15fsxcacfxde} auf den Betrag nach
Abzug von Wachstumskapital und Exploration. Latente Steuern, Verlustvortr\"age
und L\"anderunterschiede bleiben unber\"ucksichtigt; die Aufgabenstellung
dieser Analyse ist eine grobe Flow-Rechnung und keine Steuerplanung.}

\annahme{\textbf{Endwert und Horizont.} Der Anleger h\"alt am Ende des
Horizonts eine Beteiligung, und sie ist etwas wert. Angesetzt wird der
bilanzielle Buchwert des Eigenkapitals je Aktie von
\USDje{\FlowEndwert}\quelleFS{2025}{6}\quelleMerken{s15fsxcacfxg} -- belegt, unabh\"angig vom
Einstiegskurs und f\"ur eine \"uber Buchwert notierende Aktie vorsichtig. Ein
Bewertungsvielfaches w\"are ein plausibel aussehender Platzhalter und
unterbliebe damit der Belegpflicht. Was der heutige Kurs statt dessen
voraussetzt, rechnet Abschnitt~\ref{sec:endwert} aus. Gerechnet wird \"uber
f\"unf, zehn und f\"unfzehn Jahre; ein interner Zinsfu{\ss} ohne Horizont ist
eine durch eine unausgesprochene Annahme gew\"ahlte Zahl.}

\annahme{\textbf{Vorsteuerbetrachtung auf Anlegerebene.} Die Rechnung
ber\"ucksichtigt die Ertragsteuern des Unternehmens, nicht aber die Steuern
des Anlegers und keine Quellensteuer. Die ausgewiesenen internen
Zinsf\"u{\ss}e sind deshalb Vorsteuerrenditen aus Sicht des Anlegers. Da alle
Szenarien gleich behandelt werden, bleibt die Rangfolge unber\"uhrt.}

\subsection{Finanzierungskapazit\"at des Wachstumsprogramms}
\label{sec:kapazitaet}

Die erste Frage lautet, ob das ver\"offentlichte Bauprogramm aus eigener Kraft
getragen wird. Anfangsbestand ist die Liquidit\"at zum \Datum{30.\,Juni 2026}
von \MioUSDexakt{\HJLiquiditaet}; jedes Jahr kommt die Marge nach
Gleichung~\eqref{eq:marge-je-unze} hinzu, ab gehen Wachstumskapital,
Exploration, Steuer und die auf den neuen Satz hochgerechnete Dividende.

\begin{table}[htbp]
\centering
\caption{Finanzierungskapazit\"at im mittleren Szenario (erzielter Goldpreis
2025). Betr\"age in Mio.\,USD. \emph{Kapazit\"atsrechnung, keine Vorhersage:}
Sie sagt, was das Programm an Mitteln verlangt, nicht welcher Goldpreis
eintritt. Sie endet mit dem Bauprogramm, weil eine Fortschreibung dar\"uber
hinaus Kasse ohne unterstellte Verwendung auflaufen lie{\ss}e.}
\label{tab:flowkapazitaet}
\small
\begin{tabular}{lrrrrrr}
\toprule
Jahr & Marge & Investition & Steuer & Dividende & Ver\"anderung & Bestand\\
\midrule
\tabellenkoerper{data/flow_kapazitaet}
\bottomrule
\end{tabular}
\end{table}

\FloatBarrier

Der niedrigste Bestand des Pfades betr\"agt
\MioUSDexakt{\FlowBestandTief}, der Endbestand
\MioUSDexakt{\FlowBestandEnde}. F\"allt der erzielte Preis dagegen auf den
Stand von 2024 zur\"uck, sinkt der Bestand auf
\MioUSDexakt{\FlowBestandTiefUnten} und wird damit negativ.

\bk{Das Programm ist finanzierbar, aber nicht unter jedem Preis aus eigener
Kraft. Im mittleren und im oberen Szenario tr\"agt der laufende Cash-Flow es
vollst\"andig. Im unteren reicht er nicht: Es entst\"unde eine L\"ucke, die
\"uber die unbeanspruchte Kreditlinie von
\MioUSDexakt{\FazilitaetUnbeansprucht} zu schlie{\ss}en w\"are -- verf\"ugbar,
aber eben eine Verschuldung, die das Unternehmen heute nicht hat. Das ist die
eigentliche Verwundbarkeit dieses Falls, und sie liegt nicht in der Bilanz,
sondern im Investitionsplan.}

\subsection{Eichung des Modells gegen das Jahr 2025}
\label{sec:eichung}

Ein Modell, das nicht gegen ein Ist-Jahr gehalten wird, ist eine Behauptung.
Auf das Gesch\"aftsjahr 2025 angewandt ergibt die Gleichung aus
Abschnitt~\ref{sec:annahmen}:

\begin{table}[htbp]
\centering
\caption{Eichung des Modells am Gesch\"aftsjahr 2025. Betr\"age in Mio.\,USD.}
\label{tab:floweichung}
\begin{tabular}{lr}
\toprule
Posten & Betrag\\
\midrule
Marge je Unze $\times$ verkaufte Unzen & \MioUSDexakt{\FlowEichMarge}\\
abz\"uglich Wachstumskapital & \MioUSDexakt{\FlowEichWachstum}\\
abz\"uglich laufende Ertragsteuern & \MioUSDexakt{\SteuerLaufend}\\
\midrule
Modellwert & \MioUSDexakt{\FlowEichModell}\\
Tats\"achlicher freier Cash-Flow & \MioUSDexakt{\FlowEichIst}\\
\midrule
Differenz & \MioUSDexakt{\FlowEichDifferenz}\\
davon benannte Einmalposten & \MioUSDexakt{\FlowEichErklaert}\\
\midrule
Ungekl\"arter Rest & \MioUSDexakt{\FlowEichRest}\\
\bottomrule
\end{tabular}
\end{table}

\FloatBarrier

Die benannten Einmalposten sind vier: die Ausbuchung der von Argonaut
\"ubernommenen Goldtermingesch\"afte mit \MioUSDexakt{\HedgeAusbuchung}, die
Vorauszahlung auf k\"unftige Goldlieferungen mit
\MioUSDexakt{\GoldVorauszahlung}, der Aufbau des Umlaufverm\"ogens samt
gezahlter Steuern mit \MioUSDexakt{\UmlaufvermoegenSteuern} und die
aktivierten Bauzinsen mit \MioUSDexakt{\ZinsAktiviert}.\quelleFS{2025}{9}\quelleMerken{s16fsxcacfxj}
Keiner von ihnen steht seinem Wesen nach in einer Kostenkennzahl.

\bk{Der ungekl\"arte Rest von \MioUSDexakt{\FlowEichRest} auf eine Bruttomarge
von \MioUSDexakt{\FlowEichMarge} ist eine Rundungsdifferenz. Das Modell bildet
das Ist-Jahr also ab, sobald man die Einmalposten benennt -- und es ist damit
brauchbar, solange man wei{\ss}, dass es solche Posten nicht vorhersieht.}

\subsection{Die Rechnung des Anlegers}
\label{sec:anleger}

Die zweite Frage lautet, was ein Anleger verdient, der heute zu
\USDje{\FlowEinstiegskurs} einsteigt. Die Zahlungsreihe besteht aus dem Kurs
als Auszahlung, dem verteilbaren Cash-Flow je Aktie in jedem Jahr und dem
Endwert der Beteiligung im letzten Jahr.

Der Ma{\ss}stab, an dem der interne Zinsfu{\ss} gemessen wird, ist keine
konstruierte Kapitalkostengr\"o{\ss}e, sondern eine belegte Gr\"o{\ss}e des
Unternehmens selbst: seine Eigenkapitalrendite auf das durchschnittliche
Eigenkapital, gerechnet auf das \emph{bereinigte} Ergebnis. Sie betr\"agt
\Pct{\FlowHuerde}. Auf das berichtete Ergebnis gerechnet w\"aren es
\Pct{\FlowHuerdeBerichtet}; diese Zahl tr\"agt die Wertaufholung und den
Ver\"au{\ss}erungsgewinn aus Abschnitt~\ref{sec:abschluss} und taugt nicht als
Ma{\ss}stab f\"ur eine Daueranlage.

\begin{table}[htbp]
\centering
\caption{Interner Zinsfu{\ss} der Anlegerreihe beim heutigen Kurs von
\USDje{\FlowEinstiegskurs}, in Prozent. Ma{\ss}stab: \Pct{\FlowHuerde}.}
\label{tab:flowirr}
\begin{tabular}{lrrrr}
\toprule
Szenario & Goldpreis & 5 Jahre & 10 Jahre & 15 Jahre\\
\midrule
\tabellenkoerper{data/flow_irr}
\bottomrule
\end{tabular}
\end{table}

\FloatBarrier

\bk{Kein Szenario und kein Horizont erreicht den Ma{\ss}stab. Selbst im
oberen Szenario -- dem Preis des ersten Halbjahres 2026, also dem h\"ochsten,
den dieses Unternehmen je erzielt hat, dauerhaft fortgeschrieben -- bleibt der
interne Zinsfu{\ss} \"uber f\"unfzehn Jahre bei \Pct{\FlowIRRCFuenfzehn} und
damit deutlich unter \Pct{\FlowHuerde}.}

\subsection{Sensitivit\"aten}
\label{sec:sensitiv}

Zwei Annahmen entscheiden das Ergebnis st\"arker als die \"ubrigen. Beide
werden durchgerechnet, statt sie zu behaupten.

\begin{table}[htbp]
\centering
\caption{Interner Zinsfu{\ss} im mittleren Szenario bei abweichender Dauer des
Bauprogramms, in Prozent.}
\label{tab:flowsensitiv}
\begin{tabular}{lrrr}
\toprule
Wachstumskapital & 5 Jahre & 10 Jahre & 15 Jahre\\
\midrule
\tabellenkoerper{data/flow_sensitiv}
\bottomrule
\end{tabular}
\end{table}

\FloatBarrier

Die Spanne zwischen \enquote{gar kein Wachstumskapital} und
\enquote{dauerhaft fortgef\"uhrt} betr\"agt \"uber zehn Jahre nur wenige
Prozentpunkte: \Pct{\FlowIRROhneBauZehn} gegen
\Pct{\FlowIRRDauerbauZehn}. Die Annahme, die am meisten nach Willk\"ur
aussieht, entscheidet das Ergebnis also gerade nicht.

\subsubsection{Der Endwert, der \"uber das Ergebnis entscheidet}
\label{sec:endwert}

Anders steht es mit dem Endwert. Tabelle~\ref{tab:flowendwert} kehrt die
Rechnung um und nennt den Endwert je Aktie, den der \emph{heutige} Kurs
bereits voraussetzt, damit der interne Zinsfu{\ss} den Ma{\ss}stab trifft.

\begin{table}[htbp]
\centering
\caption{Erforderlicher Endwert je Aktie in USD und als Vielfaches des
Buchwerts von \USDje{\FlowEndwert}, damit der interne Zinsfu{\ss} beim
heutigen Kurs \Pct{\FlowHuerde} erreicht.}
\label{tab:flowendwert}
\small
\begin{tabular}{lrrrrrr}
\toprule
& \multicolumn{2}{c}{5 Jahre} & \multicolumn{2}{c}{10 Jahre}
& \multicolumn{2}{c}{15 Jahre}\\
\cmidrule(lr){2-3}\cmidrule(lr){4-5}\cmidrule(lr){6-7}
Szenario & USD & Vielf. & USD & Vielf. & USD & Vielf.\\
\midrule
\tabellenkoerper{data/flow_endwert}
\bottomrule
\end{tabular}
\end{table}

\FloatBarrier

Die Gegenprobe braucht einen Ma{\ss}stab, und sie bekommt einen belegten: Das
Unternehmen selbst beziffert den Kapitalwert nach Steuern allein des
Island-Gold-Distrikts auf \MrdUSD{\IGDKapitalwert}, das sind
\USDje{\FlowIGDJeAktie} je Aktie -- bei einer Preisannahme von
\JeUnze{\IGDGoldpreis}. Tabelle~\ref{tab:flowendwertvariante} rechnet die
Einstiegsschwelle deshalb ein zweites Mal, mit diesem Wert als Endwert.

\begin{table}[htbp]
\centering
\caption{Einstiegskurs \"uber f\"unfzehn Jahre in USD bei zwei belegten
Endwerten. Aktueller Kurs: \USDje{\FlowEinstiegskurs}.}
\label{tab:flowendwertvariante}
\begin{tabular}{lrr}
\toprule
Szenario & Endwert Buchwert & Endwert Kapitalwert Island Gold\\
\midrule
\tabellenkoerper{data/flow_endwertvariante}
\bottomrule
\end{tabular}
\end{table}

\FloatBarrier

\bk{Diese Tabelle ist die Belastungsprobe des ganzen Teils. Sie ersetzt die
vorsichtigste belegbare Annahme durch die g\"unstigste belegbare -- durch eine
Angabe des Unternehmens \"uber sein bestes Projekt bei einem Goldpreis, der
\"uber allen drei Szenarien liegt. Selbst dann liegt die Schwelle im oberen
Szenario bei \USDje{\FlowSchwelleIGDBeste}, und der heutige Kurs liegt
\Pct{\FlowAbstandIGDBeste} dar\"uber. Das Ergebnis h\"angt also nicht an der
Endwertannahme.}
